% !PRINTED 418
\weberpage{418}
%<<CONT
\gesperrt{Polyëdergruppen} identisch sein muss. Wir wollen hier nur den interessantesten Fall eingehender behandeln, dass $P$ die \gesperrt{Ikosaëdergruppe} ist, die wir überdies ohne Beschränkung der Allgemeinheit in der Normalform (\S.\,58) annehmen können.

Für diese Gruppe haben wir in \S.\,60 die drei Grundformen abgeleitet. Bezeichnen wir die Variablen dieser Grundformen mit $y_1, y_2$, so sind es die drei homogenen Functionen $12^{\mathrm{ten}}$, $20^{\mathrm{sten}}$ und $30^{\mathrm{sten}}$ Grades:
\weq{9}{\begin{aligned}
f(y) &= y_1 y_2 (y_1^{10} + 11 y_1^5 y_2^5 - y_2^{10}) \\
H(y) &= -y_1^{20} - y_2^{20} + 228 (y_1^{15} y_2^5 - y_2^{15} y_1^5) - 494 y_1^{10} y_2^{10} \\
T(y) &= y_1^{30} + y_2^{30} + 522 (y_1^{25} y_2^5 - y_1^5 y_2^{25}) - 10005 (y_1^{20} y_2^{10} + y_2^{20} y_1^{10}),
\end{aligned}}
zwischen denen noch die Relation
\weq{10}{T^2 + H^3 = 1728 f^5}
besteht.

Wenn wir in dem Quotienten $T^2 : f^5$ für das Verhältniss $y_1 : y_2$ einen der 60 Werthe $\vartheta_\alpha$ setzen, die aus $\vartheta$ durch die Ikosaëdersubstitutionen entstehen, so erhält dieser Quotient immer denselben Werth, der sich also rational durch die Coëfficienten der Gleichung (8), d.\,h. rational durch $z$ ausdrücken lässt. Wir können ihn geradezu gleich $z$ setzen, und erhalten für die Resolvente (8) die Form
\weq{11}{T^2 - zf^5 = 0,}
was wegen (10) mit
\[
H^3 - (1728 - z) f^5 = 0
\]
gleichbedeutend ist. Diese Gleichung ist in Bezug auf $y_1, y_2$ homogen und vom $60^{\mathrm{sten}}$ Grade, und wenn wir $y_2 = 1$ setzen, so sind ihre 60 Wurzeln $y_1 = \vartheta_\alpha$. Es ist die \gesperrt{Ikosaëdergleichung}, die schon im \S.\,60 definirt war, und die also ausführlich, in der Form (11) geschrieben, so lautet:
\[
(\vartheta^{30} + 522 \vartheta^{25} - 10005 \vartheta^{20} - 10005 \vartheta^{10} - 522 \vartheta^5 + 1)^2 = z \vartheta^5 (\vartheta^{10} + 11 \vartheta^5 - 1)^5.
\]

Das Problem, wie wir es also jetzt gefasst haben, kommt auf die Frage hinaus:

\gesperrt{Welche Gleichungen lassen sich durch die Ikosaëdergleichung lösen?}
